\hypertarget{chauxeenes-de-caractuxe8res}{%
\section{Chaînes de caractères}\label{chauxeenes-de-caractuxe8res}}

Les chaînes de caractères (\emph{strings}) ne sont pas comme des
entiers, des nombres à virgule flottante ou des booléens. Une chaîne est
une séquence (\emph{sequence}), c\textquotesingle est-à-dire une
collection ordonnée d\textquotesingle autres valeurs. Dans ce chapitre,
vous verrez comment accéder aux caractères qui composent une chaîne et
vous apprendrez certaines des méthodes proposées par les chaînes.

\hypertarget{une-chauxeene-est-une-suxe9quence}{%
\subsection{Une chaîne est une
séquence}\label{une-chauxeene-est-une-suxe9quence}}

Une chaîne est une séquence de caractères. Vous pouvez accéder aux
caractères un à la fois avec l\textquotesingle opérateur de crochet :

\begin{Shaded}
\begin{Highlighting}[]
\OperatorTok{\textgreater{}\textgreater{}\textgreater{}}\NormalTok{ fruit }\OperatorTok{=} \StringTok{\textquotesingle{}banane\textquotesingle{}}
\OperatorTok{\textgreater{}\textgreater{}\textgreater{}}\NormalTok{ lettre }\OperatorTok{=}\NormalTok{ fruit[}\DecValTok{1}\NormalTok{]}
\end{Highlighting}
\end{Shaded}

La seconde instruction sélectionne le caractère numéro 1 de
\texttt{fruit} et l\textquotesingle affecte à \texttt{lettre}.
L\textquotesingle expression entre crochets est appelée un indice
(\emph{index}). L\textquotesingle indice indique quel caractère de la
séquence vous souhaitez (d\textquotesingle où son nom). Mais vous
pourriez ne pas obtenir ce que vous attendez :

\begin{Shaded}
\begin{Highlighting}[]
\OperatorTok{\textgreater{}\textgreater{}\textgreater{}}\NormalTok{ lettre}
\CommentTok{\textquotesingle{}a\textquotesingle{}}
\end{Highlighting}
\end{Shaded}

Pour la plupart des gens, la première lettre de
\texttt{\textquotesingle{}banane\textquotesingle{}} est \texttt{b}, pas
\texttt{a}. Mais pour les informaticiens, l\textquotesingle indice est
un décalage par rapport au début de la chaîne, et le décalage de la
première lettre est zéro.

\begin{Shaded}
\begin{Highlighting}[]
\OperatorTok{\textgreater{}\textgreater{}\textgreater{}}\NormalTok{ lettre }\OperatorTok{=}\NormalTok{ fruit[}\DecValTok{0}\NormalTok{]}
\OperatorTok{\textgreater{}\textgreater{}\textgreater{}}\NormalTok{ lettre}
\CommentTok{\textquotesingle{}b\textquotesingle{}}
\end{Highlighting}
\end{Shaded}

Ainsi, \texttt{b} est la 0-ième lettre (« zéro-ième ») de
\texttt{\textquotesingle{}banane\textquotesingle{}}, \texttt{a} est la
1-ère lettre, et \texttt{n} est la 2-ème lettre. Vous pouvez utiliser
n\textquotesingle importe quelle expression, y compris des variables et
des opérateurs, comme indice, mais la valeur de l\textquotesingle indice
doit être un entier. Sinon, vous obtenez une exception :

\begin{Shaded}
\begin{Highlighting}[]
\OperatorTok{\textgreater{}\textgreater{}\textgreater{}}\NormalTok{ lettre }\OperatorTok{=}\NormalTok{ fruit[}\FloatTok{1.5}\NormalTok{]}
\PreprocessorTok{TypeError}\NormalTok{: string indices must be integers}
\end{Highlighting}
\end{Shaded}

\hypertarget{len}{%
\subsection{\texorpdfstring{\texttt{len}}{len}}\label{len}}

\texttt{len} est une fonction intégrée qui renvoie le nombre de
caractères dans une chaîne :

\begin{Shaded}
\begin{Highlighting}[]
\OperatorTok{\textgreater{}\textgreater{}\textgreater{}}\NormalTok{ fruit }\OperatorTok{=} \StringTok{\textquotesingle{}banane\textquotesingle{}}
\OperatorTok{\textgreater{}\textgreater{}\textgreater{}} \BuiltInTok{len}\NormalTok{(fruit)}
\DecValTok{6}
\end{Highlighting}
\end{Shaded}

Pour obtenir la dernière lettre d\textquotesingle une chaîne, vous
pourriez être tenté d\textquotesingle essayer quelque chose comme ceci :

\begin{Shaded}
\begin{Highlighting}[]
\OperatorTok{\textgreater{}\textgreater{}\textgreater{}}\NormalTok{ longueur }\OperatorTok{=} \BuiltInTok{len}\NormalTok{(fruit)}
\OperatorTok{\textgreater{}\textgreater{}\textgreater{}}\NormalTok{ derniere }\OperatorTok{=}\NormalTok{ fruit[longueur]}
\PreprocessorTok{IndexError}\NormalTok{: string index out of }\BuiltInTok{range}
\end{Highlighting}
\end{Shaded}

La raison de cette erreur est qu\textquotesingle il n\textquotesingle y
a pas de lettre à l\textquotesingle indice 6 dans
\texttt{\textquotesingle{}banane\textquotesingle{}}. Les indices vont de
0 à 5. Pour obtenir le dernier caractère, vous devez soustraire 1 à la
longueur :

\begin{Shaded}
\begin{Highlighting}[]
\OperatorTok{\textgreater{}\textgreater{}\textgreater{}}\NormalTok{ derniere }\OperatorTok{=}\NormalTok{ fruit[longueur}\OperatorTok{{-}}\DecValTok{1}\NormalTok{]}
\OperatorTok{\textgreater{}\textgreater{}\textgreater{}}\NormalTok{ derniere}
\CommentTok{\textquotesingle{}e\textquotesingle{}}
\end{Highlighting}
\end{Shaded}

Alternativement, vous pouvez utiliser des indices négatifs, qui comptent
à rebours à partir de la fin de la chaîne. L\textquotesingle expression
\texttt{fruit{[}-1{]}} donne la dernière lettre, \texttt{fruit{[}-2{]}}
donne l\textquotesingle avant-dernière, et ainsi de suite.

\hypertarget{parcours-avec-une-boucle-for}{%
\subsection{\texorpdfstring{Parcours avec une boucle
\texttt{for}}{Parcours avec une boucle for}}\label{parcours-avec-une-boucle-for}}

Une grande partie des calculs implique le traitement
d\textquotesingle une chaîne un caractère à la fois. Souvent, ils
commencent au début, sélectionnent chaque caractère tour à tour, font
quelque chose avec, et continuent jusqu\textquotesingle à la fin. Ce
modèle de traitement est appelé un parcours (\emph{traversal}). Une
façon d\textquotesingle écrire un parcours est
d\textquotesingle utiliser une boucle \texttt{while} :

\begin{Shaded}
\begin{Highlighting}[]
\NormalTok{indice }\OperatorTok{=} \DecValTok{0}
\ControlFlowTok{while}\NormalTok{ indice }\OperatorTok{\textless{}} \BuiltInTok{len}\NormalTok{(fruit):}
\NormalTok{    lettre }\OperatorTok{=}\NormalTok{ fruit[indice]}
    \BuiltInTok{print}\NormalTok{(lettre)}
\NormalTok{    indice }\OperatorTok{=}\NormalTok{ indice }\OperatorTok{+} \DecValTok{1}
\end{Highlighting}
\end{Shaded}

Cette boucle parcourt la chaîne et affiche chaque lettre sur une ligne à
part. La condition de la boucle est
\texttt{indice\ \textless{}\ len(fruit)}, donc lorsque \texttt{indice}
est égal à la longueur de la chaîne, la condition est fausse, et le
corps de la boucle ne s\textquotesingle exécute pas.

Une autre façon d\textquotesingle écrire un parcours est
d\textquotesingle utiliser une boucle \texttt{for} :

\begin{Shaded}
\begin{Highlighting}[]
\ControlFlowTok{for}\NormalTok{ lettre }\KeywordTok{in}\NormalTok{ fruit:}
    \BuiltInTok{print}\NormalTok{(lettre)}
\end{Highlighting}
\end{Shaded}

À chaque itération de la boucle, le caractère suivant de la chaîne est
affecté à la variable \texttt{lettre}. La boucle continue
jusqu\textquotesingle à ce qu\textquotesingle il ne reste plus de
caractères.

\hypertarget{tranches-de-chauxeenes}{%
\subsection{Tranches de chaînes}\label{tranches-de-chauxeenes}}

Un segment d\textquotesingle une chaîne s\textquotesingle appelle une
tranche (\emph{slice}). La sélection d\textquotesingle une tranche est
similaire à la sélection d\textquotesingle un caractère :

\begin{Shaded}
\begin{Highlighting}[]
\OperatorTok{\textgreater{}\textgreater{}\textgreater{}}\NormalTok{ s }\OperatorTok{=} \StringTok{\textquotesingle{}Monty Python\textquotesingle{}}
\OperatorTok{\textgreater{}\textgreater{}\textgreater{}}\NormalTok{ s[}\DecValTok{0}\NormalTok{:}\DecValTok{5}\NormalTok{]}
\CommentTok{\textquotesingle{}Monty\textquotesingle{}}
\OperatorTok{\textgreater{}\textgreater{}\textgreater{}}\NormalTok{ s[}\DecValTok{6}\NormalTok{:}\DecValTok{12}\NormalTok{]}
\CommentTok{\textquotesingle{}Python\textquotesingle{}}
\end{Highlighting}
\end{Shaded}

L\textquotesingle opérateur \texttt{{[}n:m{]}} renvoie la partie de la
chaîne du n-ième caractère au m-ième caractère, en incluant le premier
mais en excluant le dernier.

Si vous omettez le premier indice, la tranche commence au début de la
chaîne. Si vous omettez le second, la tranche va jusqu\textquotesingle à
la fin de la chaîne :

\begin{Shaded}
\begin{Highlighting}[]
\OperatorTok{\textgreater{}\textgreater{}\textgreater{}}\NormalTok{ fruit }\OperatorTok{=} \StringTok{\textquotesingle{}banane\textquotesingle{}}
\OperatorTok{\textgreater{}\textgreater{}\textgreater{}}\NormalTok{ fruit[:}\DecValTok{3}\NormalTok{]}
\CommentTok{\textquotesingle{}ban\textquotesingle{}}
\OperatorTok{\textgreater{}\textgreater{}\textgreater{}}\NormalTok{ fruit[}\DecValTok{3}\NormalTok{:]}
\CommentTok{\textquotesingle{}ane\textquotesingle{}}
\end{Highlighting}
\end{Shaded}

Si le premier indice est supérieur ou égal au second, le résultat est
une chaîne vide, représentée par deux guillemets :

\begin{Shaded}
\begin{Highlighting}[]
\OperatorTok{\textgreater{}\textgreater{}\textgreater{}}\NormalTok{ fruit[}\DecValTok{3}\NormalTok{:}\DecValTok{3}\NormalTok{]}
\CommentTok{\textquotesingle{}\textquotesingle{}}
\end{Highlighting}
\end{Shaded}

Une chaîne vide ne contient aucun caractère et a une longueur de 0, mais
à part cela, c\textquotesingle est une chaîne comme une autre.

\hypertarget{les-chauxeenes-sont-immuables}{%
\subsection{Les chaînes sont
immuables}\label{les-chauxeenes-sont-immuables}}

Il est tentant d\textquotesingle utiliser l\textquotesingle opérateur de
crochet du côté gauche d\textquotesingle une affectation, avec
l\textquotesingle intention de changer un caractère dans une chaîne. Par
exemple :

\begin{Shaded}
\begin{Highlighting}[]
\OperatorTok{\textgreater{}\textgreater{}\textgreater{}}\NormalTok{ salut }\OperatorTok{=} \StringTok{\textquotesingle{}Bonjour, monde !\textquotesingle{}}
\OperatorTok{\textgreater{}\textgreater{}\textgreater{}}\NormalTok{ salut[}\DecValTok{0}\NormalTok{] }\OperatorTok{=} \StringTok{\textquotesingle{}J\textquotesingle{}}
\PreprocessorTok{TypeError}\NormalTok{: }\StringTok{\textquotesingle{}str\textquotesingle{}} \BuiltInTok{object}\NormalTok{ does }\KeywordTok{not}\NormalTok{ support item assignment}
\end{Highlighting}
\end{Shaded}

L\textquotesingle erreur indique que l\textquotesingle objet de type
chaîne ne prend pas en charge l\textquotesingle affectation
d\textquotesingle éléments. La raison en est que les chaînes sont
immuables (\emph{immutable}), ce qui signifie que vous ne pouvez pas
modifier une chaîne existante. Le mieux que vous puissiez faire est de
créer une nouvelle chaîne qui est une variante de
l\textquotesingle originale :

\begin{Shaded}
\begin{Highlighting}[]
\OperatorTok{\textgreater{}\textgreater{}\textgreater{}}\NormalTok{ salut }\OperatorTok{=} \StringTok{\textquotesingle{}Bonjour, monde !\textquotesingle{}}
\OperatorTok{\textgreater{}\textgreater{}\textgreater{}}\NormalTok{ nouveau\_salut }\OperatorTok{=} \StringTok{\textquotesingle{}J\textquotesingle{}} \OperatorTok{+}\NormalTok{ salut[}\DecValTok{1}\NormalTok{:]}
\OperatorTok{\textgreater{}\textgreater{}\textgreater{}}\NormalTok{ nouveau\_salut}
\CommentTok{\textquotesingle{}Jonjour, monde !\textquotesingle{}}
\end{Highlighting}
\end{Shaded}

\hypertarget{recherche}{%
\subsection{Recherche}\label{recherche}}

Que fait la fonction suivante ?

\begin{Shaded}
\begin{Highlighting}[]
\KeywordTok{def}\NormalTok{ trouver(mot, lettre):}
\NormalTok{    indice }\OperatorTok{=} \DecValTok{0}
    \ControlFlowTok{while}\NormalTok{ indice }\OperatorTok{\textless{}} \BuiltInTok{len}\NormalTok{(mot):}
        \ControlFlowTok{if}\NormalTok{ mot[indice] }\OperatorTok{==}\NormalTok{ lettre:}
            \ControlFlowTok{return}\NormalTok{ indice}
\NormalTok{        indice }\OperatorTok{=}\NormalTok{ indice }\OperatorTok{+} \DecValTok{1}
    \ControlFlowTok{return} \OperatorTok{{-}}\DecValTok{1}
\end{Highlighting}
\end{Shaded}

Dans un sens, \texttt{trouver} est l\textquotesingle inverse de
l\textquotesingle opérateur de crochet. Au lieu de prendre un indice et
d\textquotesingle extraire le caractère correspondant, elle prend un
caractère et trouve l\textquotesingle indice où ce caractère apparaît.
Si le caractère n\textquotesingle est pas trouvé, la fonction renvoie
\texttt{-1}.

C\textquotesingle est le premier exemple que nous voyons
d\textquotesingle une instruction \texttt{return} à
l\textquotesingle intérieur d\textquotesingle une boucle. Si
\texttt{mot{[}indice{]}\ ==\ lettre}, la fonction sort immédiatement de
la boucle et retourne le résultat. Si le caractère
n\textquotesingle apparaît pas dans la chaîne, le programme parcourt
normalement toute la boucle et retourne \texttt{-1}.

Ce modèle de calcul, qui traverse une séquence et retourne un résultat
lorsqu\textquotesingle il trouve ce qu\textquotesingle il cherche, est
appelé une recherche (\emph{search}).

\hypertarget{boucles-et-comptage}{%
\subsection{Boucles et comptage}\label{boucles-et-comptage}}

Le programme suivant compte le nombre de fois que la lettre \texttt{a}
apparaît dans une chaîne :

\begin{Shaded}
\begin{Highlighting}[]
\NormalTok{mot }\OperatorTok{=} \StringTok{\textquotesingle{}banane\textquotesingle{}}
\NormalTok{compteur }\OperatorTok{=} \DecValTok{0}
\ControlFlowTok{for}\NormalTok{ lettre }\KeywordTok{in}\NormalTok{ mot:}
    \ControlFlowTok{if}\NormalTok{ lettre }\OperatorTok{==} \StringTok{\textquotesingle{}a\textquotesingle{}}\NormalTok{:}
\NormalTok{        compteur }\OperatorTok{=}\NormalTok{ compteur }\OperatorTok{+} \DecValTok{1}
\BuiltInTok{print}\NormalTok{(compteur)}
\end{Highlighting}
\end{Shaded}

Ce programme démontre un autre modèle de calcul appelé un compteur
(\emph{counter}). La variable \texttt{compteur} est initialisée à 0 et
est incrémentée à chaque fois qu\textquotesingle un \texttt{a} est
trouvé. À la sortie de la boucle, \texttt{compteur} contient le résultat
final.

\hypertarget{muxe9thodes-de-chauxeenes-de-caractuxe8res}{%
\subsection{Méthodes de chaînes de
caractères}\label{muxe9thodes-de-chauxeenes-de-caractuxe8res}}

Les chaînes de caractères fournissent des méthodes qui effectuent une
variété d\textquotesingle opérations utiles. Une méthode (\emph{method})
est similaire à une fonction : elle prend des arguments et retourne une
valeur, mais la syntaxe est différente.

Par exemple, la méthode \texttt{upper} prend une chaîne et retourne une
nouvelle chaîne avec toutes les lettres en majuscules. Au lieu de la
syntaxe de fonction \texttt{upper(mot)}, elle utilise la syntaxe de
méthode \texttt{mot.upper()} :

\begin{Shaded}
\begin{Highlighting}[]
\OperatorTok{\textgreater{}\textgreater{}\textgreater{}}\NormalTok{ mot }\OperatorTok{=} \StringTok{\textquotesingle{}banane\textquotesingle{}}
\OperatorTok{\textgreater{}\textgreater{}\textgreater{}}\NormalTok{ nouveau\_mot }\OperatorTok{=}\NormalTok{ mot.upper()}
\OperatorTok{\textgreater{}\textgreater{}\textgreater{}}\NormalTok{ nouveau\_mot}
\CommentTok{\textquotesingle{}BANANE\textquotesingle{}}
\end{Highlighting}
\end{Shaded}

Cette forme de notation avec un point indique le nom de la méthode,
\texttt{upper}, et le nom de la chaîne sur laquelle appliquer la
méthode, \texttt{mot}. Un appel de méthode s\textquotesingle appelle une
invocation (\emph{invocation}) ; dans ce cas, nous pourrions dire que
nous invoquons \texttt{upper} sur \texttt{mot}.

En fait, il existe une méthode de chaîne nommée \texttt{find} qui est
remarquablement similaire à la fonction \texttt{trouver} que nous avons
écrite :

\begin{Shaded}
\begin{Highlighting}[]
\OperatorTok{\textgreater{}\textgreater{}\textgreater{}}\NormalTok{ mot }\OperatorTok{=} \StringTok{\textquotesingle{}banane\textquotesingle{}}
\OperatorTok{\textgreater{}\textgreater{}\textgreater{}}\NormalTok{ index }\OperatorTok{=}\NormalTok{ mot.find(}\StringTok{\textquotesingle{}a\textquotesingle{}}\NormalTok{)}
\OperatorTok{\textgreater{}\textgreater{}\textgreater{}}\NormalTok{ index}
\DecValTok{1}
\end{Highlighting}
\end{Shaded}

\hypertarget{lopuxe9rateur-in}{%
\subsection{\texorpdfstring{L\textquotesingle opérateur
\texttt{in}}{L\textquotesingle opérateur in}}\label{lopuxe9rateur-in}}

Le mot-clé \texttt{in} est un opérateur booléen qui prend deux chaînes
et renvoie \texttt{True} si la première apparaît en tant que sous-chaîne
dans la seconde :

\begin{Shaded}
\begin{Highlighting}[]
\OperatorTok{\textgreater{}\textgreater{}\textgreater{}} \StringTok{\textquotesingle{}a\textquotesingle{}} \KeywordTok{in} \StringTok{\textquotesingle{}banane\textquotesingle{}}
\VariableTok{True}
\OperatorTok{\textgreater{}\textgreater{}\textgreater{}} \StringTok{\textquotesingle{}grain\textquotesingle{}} \KeywordTok{in} \StringTok{\textquotesingle{}banane\textquotesingle{}}
\VariableTok{False}
\end{Highlighting}
\end{Shaded}

\hypertarget{comparaison-de-chauxeenes}{%
\subsection{Comparaison de chaînes}\label{comparaison-de-chauxeenes}}

L\textquotesingle opérateur relationnel fonctionne sur les chaînes de
caractères. Pour voir si deux chaînes sont égales :

\begin{Shaded}
\begin{Highlighting}[]
\ControlFlowTok{if}\NormalTok{ mot }\OperatorTok{==} \StringTok{\textquotesingle{}banane\textquotesingle{}}\NormalTok{:}
    \BuiltInTok{print}\NormalTok{(}\StringTok{"D\textquotesingle{}accord, des bananes."}\NormalTok{)}
\end{Highlighting}
\end{Shaded}

Les autres opérations relationnelles sont utiles pour classer les mots
par ordre alphabétique :

\begin{Shaded}
\begin{Highlighting}[]
\ControlFlowTok{if}\NormalTok{ mot }\OperatorTok{\textless{}} \StringTok{\textquotesingle{}banane\textquotesingle{}}\NormalTok{:}
    \BuiltInTok{print}\NormalTok{(}\StringTok{\textquotesingle{}Votre mot, \textquotesingle{}} \OperatorTok{+}\NormalTok{ mot }\OperatorTok{+} \StringTok{\textquotesingle{}, vient avant banane.\textquotesingle{}}\NormalTok{)}
\ControlFlowTok{elif}\NormalTok{ mot }\OperatorTok{\textgreater{}} \StringTok{\textquotesingle{}banane\textquotesingle{}}\NormalTok{:}
    \BuiltInTok{print}\NormalTok{(}\StringTok{\textquotesingle{}Votre mot, \textquotesingle{}} \OperatorTok{+}\NormalTok{ mot }\OperatorTok{+} \StringTok{\textquotesingle{}, vient après banane.\textquotesingle{}}\NormalTok{)}
\ControlFlowTok{else}\NormalTok{:}
    \BuiltInTok{print}\NormalTok{(}\StringTok{"D\textquotesingle{}accord, des bananes."}\NormalTok{)}
\end{Highlighting}
\end{Shaded}

Python ne gère pas les majuscules et les minuscules de la même façon que
les humains. Toutes les lettres majuscules viennent avant toutes les
lettres minuscules, donc :

\begin{Shaded}
\begin{Highlighting}[]
\NormalTok{Votre mot, Ananas, vient avant banane.}
\end{Highlighting}
\end{Shaded}

\hypertarget{duxe9bogage}{%
\subsection{Débogage}\label{duxe9bogage}}

Lorsque vous utilisez des indices pour parcourir les valeurs
d\textquotesingle une séquence, il est facile de se tromper sur le début
et la fin de l\textquotesingle itération. Voici une fonction qui est
censée comparer deux mots et renvoyer \texttt{True} si
l\textquotesingle un d\textquotesingle eux est l\textquotesingle inverse
de l\textquotesingle autre, mais elle contient deux erreurs :

\begin{Shaded}
\begin{Highlighting}[]
\KeywordTok{def}\NormalTok{ est\_inverse(mot1, mot2):}
    \ControlFlowTok{if} \BuiltInTok{len}\NormalTok{(mot1) }\OperatorTok{!=} \BuiltInTok{len}\NormalTok{(mot2):}
        \ControlFlowTok{return} \VariableTok{False}

\NormalTok{    i }\OperatorTok{=} \DecValTok{0}
\NormalTok{    j }\OperatorTok{=} \BuiltInTok{len}\NormalTok{(mot2)}

    \ControlFlowTok{while}\NormalTok{ j }\OperatorTok{\textgreater{}} \DecValTok{0}\NormalTok{:}
        \ControlFlowTok{if}\NormalTok{ mot1[i] }\OperatorTok{!=}\NormalTok{ mot2[j]:}
            \ControlFlowTok{return} \VariableTok{False}
\NormalTok{        i }\OperatorTok{=}\NormalTok{ i }\OperatorTok{+} \DecValTok{1}
\NormalTok{        j }\OperatorTok{=}\NormalTok{ j }\OperatorTok{{-}} \DecValTok{1}

    \ControlFlowTok{return} \VariableTok{True}
\end{Highlighting}
\end{Shaded}

Si nous testons cette fonction avec les mots \texttt{"pots"} et
\texttt{"stop"}, nous nous attendons à ce qu\textquotesingle elle
renvoie \texttt{True}, mais elle provoque une erreur \texttt{IndexError}
:

\begin{Shaded}
\begin{Highlighting}[]
\OperatorTok{\textgreater{}\textgreater{}\textgreater{}}\NormalTok{ est\_inverse(}\StringTok{\textquotesingle{}pots\textquotesingle{}}\NormalTok{, }\StringTok{\textquotesingle{}stop\textquotesingle{}}\NormalTok{)}
\PreprocessorTok{IndexError}\NormalTok{: string index out of }\BuiltInTok{range}
\end{Highlighting}
\end{Shaded}

Pour le débogage, il est judicieux d\textquotesingle insérer une
instruction \texttt{print(i,\ j)} avant la ligne fautive pour afficher
les indices. Dans ce cas, lors du premier passage dans la boucle, la
valeur de \texttt{j} est 4, ce qui est hors limites pour la chaîne
\texttt{\textquotesingle{}stop\textquotesingle{}}. La bonne valeur de
départ est \texttt{len(mot2)\ -\ 1}.

Si vous corrigez cela et lancez le programme à nouveau, il ne produit
plus d\textquotesingle erreur, mais il donne un résultat incorrect. Je
vous laisse trouver et corriger la deuxième erreur !

\hypertarget{glossaire}{%
\subsection{Glossaire}\label{glossaire}}

séquence sequence Un ensemble ordonné, c\textquotesingle est-à-dire une
collection de valeurs où chaque valeur est identifiée par un indice
entier.

élément item L\textquotesingle une des valeurs composant une séquence.

indice index Une valeur entière utilisée pour sélectionner un élément
d\textquotesingle une séquence, comme un caractère dans une chaîne.

tranche slice Une partie d\textquotesingle une chaîne spécifiée par une
plage d\textquotesingle indices.

parcours traversal Itérer à travers les éléments d\textquotesingle une
séquence, en effectuant une opération similaire sur chacun
d\textquotesingle eux.

immuable immutable La propriété d\textquotesingle une séquence dont les
éléments ne peuvent pas être modifiés.

recherche search Un modèle de parcours qui s\textquotesingle arrête
lorsqu\textquotesingle il trouve ce qu\textquotesingle il cherche.

compteur counter Une variable utilisée pour compter quelque chose,
souvent initialisée à zéro puis incrémentée.

invocation invocation Une instruction qui appelle une méthode.

méthode method Une fonction associée à un objet et appelée en utilisant
la notation pointée.

\hypertarget{exercices}{%
\subsection{Exercices}\label{exercices}}

\textbf{Exercice 1}

Lisez la documentation des méthodes de chaînes de caractères à
l\textquotesingle adresse
\url{https://docs.python.org/3/library/stdtypes.html\#string-methods}.
Expérimentez avec quelques-unes d\textquotesingle entre elles pour vous
assurer de bien comprendre comment elles fonctionnent. \texttt{strip} et
\texttt{replace} sont particulièrement utiles.

La documentation utilise une syntaxe qui pourrait être source de
confusion. Par exemple, dans
\texttt{find(sub{[},\ start{[},\ end{]}{]})}, les crochets indiquent des
arguments optionnels. Ainsi, \texttt{sub} est obligatoire, mais
\texttt{start} est facultatif, et si vous incluez \texttt{start}, alors
\texttt{end} est également facultatif.

\textbf{Exercice 2}

Une tranche de chaîne peut accepter un troisième indice pour spécifier
un "pas" (\emph{step}), c\textquotesingle est-à-dire
l\textquotesingle espacement entre les caractères consécutifs. Un pas de
2 signifie tous les deux caractères, et un pas de 3 signifie tous les
trois caractères, etc.

\begin{Shaded}
\begin{Highlighting}[]
\OperatorTok{\textgreater{}\textgreater{}\textgreater{}}\NormalTok{ fruit }\OperatorTok{=} \StringTok{\textquotesingle{}banane\textquotesingle{}}
\OperatorTok{\textgreater{}\textgreater{}\textgreater{}}\NormalTok{ fruit[}\DecValTok{0}\NormalTok{:}\DecValTok{5}\NormalTok{:}\DecValTok{2}\NormalTok{]}
\CommentTok{\textquotesingle{}bnn\textquotesingle{}}
\end{Highlighting}
\end{Shaded}

Un pas de -1 fait parcourir le mot à l\textquotesingle envers, de sorte
que la tranche \texttt{{[}::-1{]}} génère une chaîne inversée.

Utilisez cette fonctionnalité pour écrire une version raccourcie de
\texttt{est\_palindrome} (une fonction qui vérifie si un mot se lit de
la même manière à l\textquotesingle endroit et à
l\textquotesingle envers).

\textbf{Exercice 3}

Les méthodes de chaînes suivantes sont toutes utiles : \texttt{islower},
\texttt{isupper}, et \texttt{isalpha}.

Écrivez une fonction nommée \texttt{lettre\_minuscule\_presente} qui
prend une chaîne de caractères comme argument et qui retourne
\texttt{True} si la chaîne contient au moins une lettre minuscule, et
\texttt{False} sinon.

\textbf{Exercice 4}

Le chiffrement de César (ROT13) consiste à remplacer chaque lettre par
la lettre se trouvant un certain nombre de places plus loin dans
l\textquotesingle alphabet.

Écrivez une fonction nommée \texttt{tourner\_mot} qui prend une chaîne
de caractères et un entier comme paramètres, et retourne une nouvelle
chaîne contenant les lettres de la chaîne d\textquotesingle origine
tournées du nombre de places spécifié par l\textquotesingle entier.

\emph{Indice : vous pourriez trouver les fonctions intégrées}
\texttt{ord} \emph{et} \texttt{chr} \emph{utiles pour convertir des
caractères en valeurs numériques et inversement.}
